% Lecture notes: numbered definitions, theorems and examples, with a table
% of contents.
\documentclass[11pt,a4paper]{article}

\usepackage[T1]{fontenc}
\usepackage[margin=2.5cm]{geometry}
\usepackage{amsmath,amssymb,amsthm}
\usepackage{graphicx}
\usepackage[hidelinks]{hyperref}

% Theorem-like environments share one counter per section.
\theoremstyle{plain}
\newtheorem{theorem}{Theorem}[section]
\newtheorem{lemma}[theorem]{Lemma}
\newtheorem{corollary}[theorem]{Corollary}
\theoremstyle{definition}
\newtheorem{definition}[theorem]{Definition}
\newtheorem{example}[theorem]{Example}
\theoremstyle{remark}
\newtheorem*{remark}{Remark}

% Shortcuts you will type a hundred times.
\newcommand{\R}{\mathbb{R}}
\newcommand{\N}{\mathbb{N}}
\newcommand{\abs}[1]{\left\lvert #1 \right\rvert}

\title{Course Name: Lecture Notes}
\author{Lecturer: Professor Name \\ Notes by Your Name}
\date{Semester, Year}

\begin{document}

\maketitle
\tableofcontents

\section{Sequences and Limits}

\begin{definition}
A sequence $(a_n)$ in $\R$ \emph{converges} to $L$ if for every
$\varepsilon > 0$ there is an $N \in \N$ such that $\abs{a_n - L} < \varepsilon$
for all $n \ge N$.
\end{definition}

\begin{example}
The sequence $a_n = 1/n$ converges to $0$: given $\varepsilon > 0$, take
$N > 1/\varepsilon$.
\end{example}

\begin{theorem}[Uniqueness of limits]\label{thm:unique}
A convergent sequence has exactly one limit.
\end{theorem}

\begin{proof}
Suppose $a_n \to L$ and $a_n \to M$. For any $\varepsilon > 0$ and $n$ large,
\[
  \abs{L - M} \le \abs{L - a_n} + \abs{a_n - M} < 2\varepsilon,
\]
so $L = M$.
\end{proof}

\begin{remark}
Theorem~\ref{thm:unique} is what makes the notation $\lim a_n$ meaningful.
\end{remark}

\section{Continuity}

\begin{definition}
A function $f \colon \R \to \R$ is \emph{continuous} at $c$ if
$\lim_{x \to c} f(x) = f(c)$.
\end{definition}

\end{document}
